% conclusion.tex

\section{Conclusion}

We presented \textsc{Instrumental}, a system for automatic synthesizer parameter recovery from audio. Given a target recording, our approach recovers the 28 parameters of a differentiable subtractive synthesizer---spanning oscillator mix, filter characteristics, ADSR envelopes, unison, reverb, and parametric EQ---by combining spectral-analysis initialization with CMA-ES evolutionary optimization under a composite spectral loss. The best configuration achieves a matching loss of 2.09, placing it in the top 0.13\% of the explored loss landscape.

A systematic evaluation of eight hypotheses yielded actionable insights for the inverse synthesis problem. Most notably, spectral shaping via parametric EQ proved more impactful than alternative waveform generators (Chebyshev waveshaping, FM synthesis), alternative loss functions (TFS envelope, CDPAM, CLAP embeddings), or alternative optimization strategies (SPSA fine-tuning, multi-start CMA-ES). This finding suggests that for subtractive synthesis targets, the ability to selectively boost frequency bands compensates for limitations in oscillator topology more effectively than increasing waveform complexity.

Our work addresses a complementary gap in the audio analysis pipeline: whereas existing tools such as pitch trackers and transcription systems extract \emph{what notes} are played, \textsc{Instrumental} extracts \emph{how the instrument sounds}---recovering the synthesizer configuration that produced a given timbre. This distinction opens practical applications in sound design, preset recovery, and audio production workflows where recreating a specific timbre is the goal rather than transcribing pitch content.

Future directions include extending the synthesizer model to support FM and wavetable architectures, incorporating learned perceptual losses trained on human similarity judgements, scaling to polyphonic audio via source separation, and integrating the recovered parameters with commercial synthesizer formats for direct preset export.
